%% kapa: exam=2 round=27 year=2016 period=1 subject=appraisal_practice subject_ko=감정평가실무 seq=4 number=4 points=10
%% source: Q-Net 감정평가사 기출문제 (cst00309) artlSeq=5185401 file=extracted/1008/5185401/제27회 감정평가사 2차시험 기출문제/제27회 감정평가사 2차 시험문제지/(1교시 시험문제지)감정평가실무.pdf
%% conversion: kordoc 4.24.51 -> tools/gen2.py (auto); review=proofread-gpt-sol; flags=none
\begin{longquestion}{【문제 4】}
감정평가사 甲은 S시장으로부터 ‘도시계획시설도로 개설사업’에 일부 편입되는 소규모 봉제공장에 대하여 ‘영업의 휴업손실’에 대한 보상평가를 의뢰받았다. 한편, 이 공장은 관계 법령에 의하면, 일부 편입에 따른 잔여시설 보수 후 재사용이 가능한 경우에 해당되는 바, 제시된 자료를 참조하여 영업장소를 이전하는 경우와 비교하여 보상평가액을 산정하시오. (10점) \par
<자료 1> 공통사항 \par
1. 사업인정고시 의제일 : 2015년 7월 1일 \par
2. 본 공장의 가동일 : 2010년 7월 1일 (적법한 허가취득) \par
3. 가격시점 : 2016년 7월 1일 \par
4. 사업시행자 : S시장 \par
5. 휴업기간 및 잔여시설 보수기간 : 4개월 \par
<자료 2> 조사내용 \par
1. 연간 영업이익 : 60,000,000원 \par
2. 영업장소 이전 후 발생하는 영업이익 감소액 : 연간 영업이익의 10 \% \par
3. 월간 고정적 비용항목 : 2,000,000원 \par
4. 전체 이전비 및 감손상당액 : 4,000,000원 \par
5. 개업비 등의 부대비용 : 1,000,000원 \par
6. 해당시설의 보수비 : 18,000,000원 \par
7. 영업규모 축소에 따른 고정자산 등의 매각손실액 : 5,000,000원 \par
\end{longquestion}
